\scanpage{30}
% archive sheet number: 29
% L1--3, continuation
est ainsi ramené à prouver que
\[
\operatorname{Hom}(Q,Q)\longrightarrow\operatorname{Hom}({\unsure{\varphi|\phi}}^*(Q),\varphi^*(Q))
\]
est bijectif. Or les deux s’identifient {\unsure{resp!|resp.}} à
% L4--5
\[
\operatorname{Hom}_{\substack{\text{grps}\\\text{sur }{\unsure{Y|S}}}}({\unsure{\mathscr P|\mathcal P}}_S,G)
\quad\text{et}\quad
\deleted{{\unsure{E|H}}}\operatorname{Hom}_Y(X-S,G)
\]
donc sont isomorphes.

% L6--9
En particulier

\noindent\vrule\hspace{.6em}\begin{minipage}{.93\linewidth}
\uline{Corollaire 3} Si $G$ est un schéma en groupes commu. de type fini sur $Y$, on a un hom. canonique \uuline{injectif}
\end{minipage}\par
\[
\operatorname{Ext}_{\substack{\text{schémas en groupes}\\\text{sur }{\unsure{Y|S}}}}(\mathscr P_{S,X/S},G)\longrightarrow H^1(X-S,G)
\]
% margin note written slanted
\margin{{\unsure{préciser la signification}}{\unsure{...}}}
% L10--15
Il \uline{semble} que c’est \uline{{\unsure{bijectif|surjectif}}} chaque fois que $S\longrightarrow Y$ est \uline{surjectif}, i.e. que les fibres de $X/Y$ sont non propres, donc \uline{affines}. Dans le cas où $S=\varnothing$, $G=G_m$, on ne {\unsure{sera|serait}} déjà pas surjectif.
